\icode[magenta]{E-MCH-MU-L4-261025-001}
（多选）
如图所示，一质点在光滑水平桌面上受水平恒力作用，先后经过\(a\)、\(b\)两点，速度方向偏转\qty{90}{\degree}。已知经过\(a\)点的速度大小为\(v\)、方向与\(ab\)连线夹角为\qty{60}{\degree}，\(ab\)连线长度为\(d\)。对质点从\(a\)到\(b\)的运动过程，下列说法正确的是\paren[]

\begin{tikzpicture}[>=stealth,font=\scriptsize]
    \draw[dashed,cyan] (0,0)--(60:2)--(4,0)node[above,black]{\(b\)} (-.5,0)--(4.5,0);
	\draw[->,red] (0,0)node[below,black]{\(a\)}--(60:.8)node[left,black]{\(v\)};
	\draw[->,blue] (4,0)--++(-30:{.5*sqrt(3)});
	\draw (0.2,0) arc(0:60:.2);
	\node at(.6,.2){\qty{60}{\degree}};
\end{tikzpicture}
\begin{choices}
	\item 最小速度为\(\frac{v}{2}\)
	\item 运动时间为\(\frac{d}{2v}\)
	\item 经过\(b\)点的速度为\(\sqrt{3}v\)
	\item 恒力方向与\(ab\)连线的夹角为\qty{45}{\degree}
\end{choices}
\begin{answer}
	BC
\end{answer}